\documentclass[10pt,a4paper]{amsart}
\usepackage[margin=19mm]{geometry}
\usepackage{amsmath,amssymb,mathtools}
\usepackage[scr=rsfs,cal=euler]{mathalfa}
\usepackage{xfrac}
\usepackage[unicode,hidelinks,bookmarks=false]{hyperref}
\newcommand{\ie}{i.e.\ }
\newcommand{\cf}{cf.\ }
\newcommand{\resp}{respectively\ }
\newcommand{\Subsection}{\subsection}
\providecommand{\nobreakdash}{\nobreak\hbox{-}}
\setlength{\emergencystretch}{2em}
\multlinegap0pt
\usepackage{bm}
\usepackage{bbm}
\usepackage{dsfont}
\usepackage{xspace}
\usepackage{cancel}
\usepackage{mathtools}
\usepackage{amsthm}
\makeatletter
\@ifundefined{theorem}{\newtheorem{theorem}{Theorem}}{}
\makeatother
\newcommand{\bG}{\mathbb G}
\newcommand{\bN}{\mathbb N}
\newcommand{\bP}{\mathbb P}
\newcommand{\cH}{\mathcal H}
\title[Multigraded Hurwitz forms]{Multigraded Hurwitz forms}
\author{Elizabeth Pratt, Luca Sodomaco, Bernd Sturmfels}
\date{Source: arXiv:2602.19563v1 (CC BY 4.0)}
\begin{document}
\maketitle

\noindent\emph{Source and licence.} Multigraded Hurwitz forms. Version arXiv:2602.19563v1 (CC BY 4.0), distributed under the Creative Commons Attribution 4.0 International licence, as recorded in the arXiv licence metadata for this exact version and shown on the abstract page https://arxiv.org/abs/2602.19563v1. Full rights notice: licenses/arxiv-2602.19563-CC-BY-4.0.txt.\par\smallskip

\noindent\textit{Editorial scope.} Counted unit: the Theorem whose source label is \texttt{thm:CIformula} in the article (source0.tex lines 1003--1014, source label \texttt{thm:CIformula}), delivered together with the article's own complete original proof of it (source0.tex lines 1016--1059, one \texttt{proof} environment of 44 lines). No mathematics is written, repaired or re-derived: statement and proof are the article's own text line for line. Required context retained uncounted: eq:multidegree (lines 417--420); thm:degree (lines 844--858); eq:degreeofhurwitz (lines 853--856). Every definition, display and label of those lines is retained unchanged. Omitted from this package and not redistributed: 1--416, 421--843, 857--1002, 1015--1015, 1060--1814. Nothing inside a retained excerpt is omitted or altered: every retained body is delivered exactly as the article's lines are written, so no omission receipt of any kind accompanies this package. Citations: the retained excerpts carry no citation command at all, so no bibliography entry of the article is transcribed and no reference is redistributed. Reference adaptation: none was needed and none was made; every reference inside the delivered unit resolves inside it, so no unresolved reference or citation is produced. Printed slips kept literal: no printed word, symbol, hypothesis or number of the counted statement or of its proof was altered. Layout and class: the article's own class is \texttt{article} with packages including standalone, amsmath, amsthm, amssymb, color, hyperref, graphicx, caption, mathtools, enumerate, verbatim, tikz, tikz-cd, tikz-3dplot; the wrapper is the lane's pinned \texttt{amsart} editorial wrapper at 10pt on A4, which restates the theorem-like environments the retained text needs and adds the article's own macro definitions that the retained text needs (\texttt{\textbackslash bG}, \texttt{\textbackslash bN}, \texttt{\textbackslash bP}, \texttt{\textbackslash cH}), with extra wrapper packages (bm, bbm, dsfont, xspace, cancel, mathtools). The delivered numbering is the unit's own. Assets: no retained span contains a figure environment, a graphics-inclusion command, a PDF-page inclusion, a table or a TikZ picture; no image file is copied into this package, the package declares no asset dependency and no image file is redistributed with it. Licence: version arXiv:2602.19563v1 of this article is distributed under the Creative Commons Attribution 4.0 International licence, as recorded in the arXiv licence metadata for this exact version and shown on the abstract page https://arxiv.org/abs/2602.19563v1. A full rights notice is delivered at licenses/arxiv-2602.19563-CC-BY-4.0.txt. Characters: every retained excerpt is pure ASCII, so the delivered engine, XeLaTeX with its default Latin Modern text font, reports no missing character.
\begin{equation}
\label{eq:multidegree}
\qquad    [X] \,\, = \,\, \sum_{|\alpha|=c} \delta_\alpha\,T^{\alpha} \qquad \in \quad A^*(\bP).
\end{equation}

\begin{theorem} \label{thm:degree}
	Let $X$ be an irreducible variety of codimension $c$ in $\bP$. 
	Fix $\alpha  \in \bN^\ell$ with $|\alpha| = c$, and
		$\delta_\alpha \geq 2$ in (\ref{eq:multidegree}),
		and $n_i \geq {\rm max}(2,1+\alpha_i)$
	 	  for $i\in[\ell]$.
Then  the multigraded Hurwitz locus  $\cH^\alpha_X$ is an irreducible variety in $\bG_\alpha$, and
%    If $\,\cH^\alpha_X$ is a hypersurface then t
the degree $u$ of the     Hurwitz form ${\rm Hu}^\alpha_X$ satisfies
\begin{equation}
\label{eq:degreeofhurwitz}
        u_i \,\,\le\,\, 2 (g_{\alpha+e_i}+\delta_\alpha - 1) \qquad \hbox{for}  \,\,\, i \in [\ell].
\end{equation}
All inequalities in (\ref{eq:degreeofhurwitz}) are equalities if the variety $X$ is polynodal.
\end{theorem}

\begin{theorem} \label{thm:CIformula}
For any vector $\alpha \in \bN^\ell$ with $|\alpha| = c$ and $\delta_\alpha \geq 2$,
the multigraded Hurwitz form ${\rm Hu}^\alpha_X$ of the generic complete intersection $X$
is a polynomial of degree $(u_1,\ldots,u_\ell)$, where
\begin{equation}
\label{eq:CIui}
     \quad u_i \,\,\,= \,\,\,2 \delta_\alpha  \, \, +\!
     \sum_{\substack{j\in[\ell]\\\alpha_j+\hat{\delta}_{ij}>0}} 
     \! \delta_{\alpha+e_i-e_j} \cdot \biggl(-\alpha_j-1-\hat{\delta}_{ij}+\sum_{p=1}^c b_{pj} \biggr)\,
\quad \hbox{for all}\,\,\, \,i \in [\ell] .
\end{equation}
\end{theorem}

\begin{proof}
The generic complete intersection $X$ is polynodal. Indeed, 
consider any  generic curve section $C = X \cap (L_1 \times \cdots \times L_\ell)$.
Then $C$ is polynodal becauses its projections into $\prod_{i \in S} \bP^{n_i}$ for $S \subseteq [\ell]$
are curves with at most nodal singularities. Therefore,  by Theorem \ref{thm:degree}, the Hurwitz form
${\rm Hu}_X^\alpha$ has degree $2 (g_{\alpha+e_i}+\delta_\alpha - 1) $
in the Pl\"ucker coordinates of the $i$th Grassmannian
$\bG(\alpha_i,n_i)$. We must show that
$2 (g_{\alpha+e_i}+\delta_\alpha - 1) $ equals
the expression for $u_i$ given in (\ref{eq:CIui}).

We shall use the adjunction formula.
Consider a curve in $\bP$ that is the intersection of
$N-1$ generic divisors of degrees $D_1,D_2,\ldots,D_{N-1}$.
The genus $g$ of this curve satisfies
\begin{equation}
\label{eq:2g-2}
  2g -2 \,\,= \,\,D_1 D_2 \cdots D_{N-1} \cdot (D_1+\cdots+D_{N-1}+K_{\bP})\,. 
  \end{equation}
Here $K_\bP$ denotes the canonical divisor of $\bP $, which is known to be
$$ K_{\bP} \,\, = \,\, \sum_{i=1}^\ell (-n_i-1) \cdot T_i . $$
 Let $D_1, \, \ldots, \, D_c$ be the divisors that define $X$, and 
 let $D_{c+1},\ldots,D_{N-1}$ be the divisors that define
 the multilinear variety $L_1 \times\cdots \times L_\ell$ which intersects $X$
  in a curve of genus $g_{\alpha+e_i}$.
  Let $n$ denote the vector $(n_1, \, \ldots, \, n_\ell)$.
  Then $D_{c+1} \cdots D_{N-1} = T^{n-\alpha}/T_i$, and therefore
      \begin{equation}\label{eq: product Di}
        D_1 \cdots D_c D_{c+1}\cdots D_{N-1} \,\,=\,\,         [X] \cdot D_{c+1} \cdots D_{N-1} \,\,
         = \sum_{\substack{j\in[\ell]\\\alpha_j+\hat{\delta}_{ij}>0}} \delta_{\alpha + e_i - e_j}T^{n - e_j}.
    \end{equation}
The rightmost factor in (\ref{eq:2g-2}) is the divisor class
\begin{equation}\label{eq: sum Di minus K}
  D_1+\cdots+D_{N-1}+K_{\bP}\, \,\,= \,\,\sum_{k=1}^\ell\left(-\alpha_k-\hat{\delta}_{ik}-1+\sum_{p=1}^c b_{pk}\right)T_k.
 \end{equation}
    We multiply \eqref{eq: product Di} with \eqref{eq: sum Di minus K}. Because we are working in the Chow ring $A^*(\bP),$ the $j$th and $k$th terms in each sum multiply to zero unless $j=k.$ The 
    resulting simplified expression is
    $$
        2g_{\alpha + e_i} -2 \,\,= \,\,
        \sum_{j=1}^\ell \delta_{\alpha + e_i - e_j} \biggl(- \alpha_j - 1 - \hat{\delta}_{ij} + \sum_{p=1}^c b_{pj} \biggr).
$$
    By adding $2 \delta_\alpha$ to this expression, we obtain the
    formula for $u_i$ that is asserted in (\ref{eq:CIui}).
\end{proof}

\end{document}
